\documentclass[11pt]{article}
\usepackage[utf8]{inputenc}
\usepackage[T1]{fontenc}
\usepackage{lmodern}
\usepackage[margin=1in]{geometry}
\usepackage{amsmath,amssymb}
\usepackage{textcomp}
\usepackage[dvipsnames]{xcolor}
\usepackage{enumitem}
\usepackage{parskip}
\usepackage{titlesec}
\usepackage[colorlinks=true,linkcolor=blue,citecolor=blue,urlcolor=blue]{hyperref}
\usepackage[numbers,sort&compress]{natbib}
\bibliographystyle{plainnat}

\titleformat{\section}{\large\bfseries}{\thesection}{1em}{}
\titleformat{\subsection}{\normalsize\bfseries}{\thesubsection}{1em}{}
\titleformat{\subsubsection}{\normalsize\bfseries\itshape}{}{0em}{}

\title{\textbf{Two Gates, Not One Score:\\ Disentangling Observability and Verifiability in LLM Bargaining Exploitation}}
\author{}
\date{}

\newcommand{\ph}{\textit{[placeholder --- fill in after experiments are run]}}

\begin{document}
\maketitle
\vspace{-2em}

\begin{quote}
\small\textit{Assembled from the six confirmed stage outputs of an \texttt{auto-research-discussion} run (seed topic B5.2), post research-taste audit (\textcolor{ForestGreen}{\textbf{PASS}}). Status: \textbf{specification / paper draft --- no experiments have been run.} All numbers below are pre-registered design targets, not results.}
\end{quote}

\begin{abstract}
Autonomous LLM agents increasingly bargain on behalf of principals who cannot observe every turn, and information asymmetry between the two negotiating agents is the default condition, not an edge case: one side may not know the other's reservation price, and either side can make claims the other has no way to check. Whether this asymmetry is a benign, even privacy-preserving, feature of a bargaining protocol or a systemic exploitation vector --- and which of its component channels a protocol designer should actually regulate --- has not been measured as a continuous, decomposed quantity. Existing work either characterizes bargaining equilibria under private valuations with no channel for a claim to be false, documents that LLM agents can deceive an unable-to-verify counterpart, or --- most recently --- compares LLM bargainers across a handful of discrete information regimes; none manipulates the two information channels as continuous, independently crossed treatments, and none adjudicates whether the resulting exploitation scales smoothly or crosses a threshold. We build a bilateral LLM bargaining testbed that independently manipulates reservation-price observability and claim verifiability, score exploitation against a ground-truth-anchored symmetric-split benchmark invariant to the manipulation itself, and separate the two channels' main effects, their relative magnitude, and the share of the verifiability effect mediated by measured deception. A pre-registered, jointly-gated statistical test then adjudicates whether the resulting relationship is a smooth dose-response curve or crosses a genuine threshold into systemic exploitation. We report which of these two pictures holds for current LLM bargainers, and which information channel is the more consequential lever for protocol designers to regulate.
\end{abstract}

\section{Introduction}

Bilateral bargaining between autonomous software agents is no longer confined to laboratory demonstrations. As large language models are delegated procurement negotiations, resale transactions, and service contracts, one condition is unavoidable: information asymmetry. A principal cannot observe every turn of the negotiation, one side may simply not know the other's reservation price, and either side can make claims --- about urgency, quality, or willingness to pay --- that the other has no way to check. Whether this default condition is dangerous is exactly the design question protocol builders face: does an agent-to-agent bargaining protocol need mandatory reservation-price disclosure, an ongoing claim-audit mechanism, or neither, because bounded asymmetry is an acceptable, even privacy-preserving, feature of a working protocol?

Existing evidence answers neither side of this question directly, and the reason is structural rather than incidental. Incomplete-information bargaining theory characterizes what happens at equilibrium once each side's valuation is private, but its formalism has no natural-language channel through which a claim could be made, let alone be false. LLM-agent work on deception and persuasion runs the opposite direction: it demonstrates that models can and do exploit a counterpart who cannot verify what they are told, but these are existence proofs of a single exploitative episode, not a measured relationship between the degree of informational advantage and the resulting surplus captured. A protocol designer needs a dose-response curve --- how much does exploitation grow as asymmetry grows, and does it grow smoothly or cross a point at which it becomes systemic rather than marginal --- and neither literature supplies one.

The nearest existing work approaches pieces of this gap without assembling them. GLEE \citep{GLEE} offers a unified, parameterized family of language-based economic games with configurable information availability, but treats information conditions as discrete settings to benchmark models under rather than a continuously swept variable, and has no channel for claim verifiability at all --- its statements are cheap talk, never audited. SOTOPIA \citep{SOTOPIA} scores rich multi-turn social scenarios along a secrecy dimension using an LLM-judge rubric, but that rubric is not anchored to the ground-truth surplus a bargaining episode actually makes available, so it cannot say how much an advantaged agent extracted. CompeteAI \citep{CompeteAI} studies long-horizon market competition with endogenous reputation feedback, which is precisely what prevents attributing an observed advantage to information structure rather than to the accumulated dynamics its own design is built to surface. And formal treatments of one-sided informational advantage --- a generalized principal-agent framework bounding how much a learning agent can be exploited \citep{LearningAgentPA}, and a Bayesian-persuasion formalization of one-directional sender-receiver games \citep{StrategicPersuasion} --- either assume away private information on one side or restrict the interaction to a single direction, neither of which matches a bilateral bargain where both sides hold private information and reputational leverage. A 2026 wave of work has since moved onto adjacent ground and sharpens, rather than closes, the gap: bilateral LLM bargaining has now been compared across discrete information regimes with honesty and credulity measured \citep{UsedCarSalesbots}, evaluator-held ground truth about a counterpart's private state is now standard diagnostic machinery \citep{TERMSBench}, and deception in an asymmetric negotiation has been monitored in real time \citep{ThinkingOutLoud} --- but the information structure in each is a fixed or categorical condition, never a pair of continuously swept, factorially crossed treatments, and none of them asks the dose-response question this paper is built around.

Turning this gap into a measurable question surfaces four distinct difficulties. Scoring how much surplus an advantaged agent captures requires a benchmark computed purely from the scenario's known parameters, since any measure that leans on what an agent discloses or an LLM judge infers would be contaminated by the very asymmetry it is supposed to quantify. Observability and verifiability are not one dial: they are mechanistically different channels --- not knowing a counterpart's value is a different obstacle from not being able to check what the counterpart says --- and collapsing them into a single asymmetry index would erase exactly the distinction a protocol designer's regulatory choice depends on. Any curve measured from LLM bargainers is entangled with which specific models and prompts were used, and no finite replication battery can fully establish that a finding is a property of information structure rather than of today's instruction-tuned models. And a claimed discontinuity in the exploitation curve is exactly the kind of artifact a coarse grid or a small sample can manufacture even when the true relationship is smooth, so a threshold claim needs more than one statistical test to survive scrutiny.

Our approach turns each of these obstacles into a specific design choice. Every exploitation reading is computed as arithmetic on evaluator-held ground truth against a symmetric-split benchmark, never on what an agent believes or discloses, so the metric cannot inherit the asymmetry it measures. Reservation-price observability and claim verifiability are implemented as two physically independent mechanisms --- a signal of controllable width, and a probabilistic audit that publicly flags false claims --- crossed factorially so a single regression can separate their individual contributions, their relative magnitude, and the share of the verifiability effect carried by a measured deception rate, rather than reporting one collapsed asymmetry score. Internal validity and cross-model generality are deliberately staged rather than conflated: one pinned model and protocol establishes the causal sweep, and a bounded replication across further model families and prompt framings measures --- without claiming to eliminate --- how much the finding depends on which agents were used. And any claim that the relationship is a genuine threshold, rather than a smooth curve, must jointly survive a model-comparison test, an effect-size criterion, and replication under a denser grid before the word threshold is used.

We evaluate this design on bilateral single-issue price bargaining, sweeping the two information-structure axes over synthetic scenarios with controlled asymmetry and extending to real negotiation listings, and comparing against an information-symmetric anchor and a scripted, memoryless opponent that rules out the possibility that any observed trend is protocol arithmetic rather than strategic use of information.

In summary, we contribute:
\begin{itemize}[leftmargin=1.4em]
\item An identification and formalization of bargaining information asymmetry as two mechanistically independent channels --- reservation-price observability and claim verifiability --- showing why collapsing them into one asymmetry index would erase the distinction a protocol designer's regulatory choice depends on.
\item A ground-truth-anchored exploitation metric, computed as arithmetic on evaluator-held reservation values and a symmetric-split benchmark, that is invariant to the information manipulation it measures.
\item A separation of the two channels' causal contributions to exploitation --- their individual main effects, their relative magnitude, and the share of the verifiability effect mediated by a measured deception rate --- from a single factorial design.
\item A pre-registered, jointly-gated statistical test for whether exploitation crosses a genuine threshold rather than scaling smoothly, together with a bounded cross-model replication that measures, without overclaiming to eliminate, how much any finding depends on the specific LLM bargainers tested.
\end{itemize}

\section{Problem Definition}

\noindent\textbf{Formal setup.} Consider a buyer $b$ and a seller $s$ bargaining over a single item via alternating offers (bounded turns; fixed message schema: offer price, accept/reject/walk-away flag, free-text claim field). Reservation values $r_b, r_s$ are drawn from controlled distributions and are known to the evaluator but private to each agent; $\mathrm{ZOPA} = \max(0, r_b - r_s)$, and analysis is restricted to positive-ZOPA episodes. One bargaining seat is fixed, by design, as the \textbf{advantaged} agent $a \in \{b, s\}$ for the episode --- an exogenous, researcher-assigned label, not one inferred post hoc from who ends up better off --- and the counterpart is the disadvantaged agent $d$. If a deal closes, it yields a realized price $p^*$. The information structure is a pair $\theta = (o, v)$: an observability level $o \in \{0,1,2,3\}$ governing the width of a signal about $r_d$ delivered to $a$ (width shrinking from unbounded at $o=0$ to zero, i.e.\ exact disclosure, at $o=3$), and a verifiability level $v \in [0,1]$, the probability with which a declared claim about private state is checked against ground truth and, if false, publicly flagged. Exploitation is $E(\theta) = (u_a(\theta) - u_a^{\text{bench}})/\mathrm{ZOPA}$, where $u_a(\theta)$ is $a$'s realized surplus and $u_a^{\text{bench}} = \mathrm{ZOPA}/2$ is the equal-split (symmetric Nash bargaining) benchmark. A condition is classified \textbf{marginal} if its mean exploitation reading is close to, but detectably above, zero --- operationalized as a mean $E$ in $(0.02, 0.10]$ --- and \textbf{systemic} if the advantaged agent's mean exploitation reading is high enough that it is capturing a large majority of the available surplus, operationalized as mean $E \geq 0.25$ with a 95\% bootstrap confidence interval that excludes $0.10$. These bands are fixed before any data are collected, precisely so that ``systemic'' is never assigned informally after seeing the results. The object of study is how $E(\theta)$ behaves as $\theta$ is swept: does it change smoothly with $o$ and $v$, and does it --- along the path of jointly maximal asymmetry, or along either axis with the other held fixed --- cross a genuine discontinuity into systemic exploitation, rather than scaling gradually through the marginal band?

\noindent\textbf{Why this is hard.} A benchmark for $E$ must be computable purely from $r_b$, $r_s$, and $\mathrm{ZOPA}$, independent of what either agent discloses, believes, or is judged to have done --- otherwise the measurement would inherit the very asymmetry it is meant to quantify. The two coordinates of $\theta$ are not one scalar: $o$ and $v$ are implemented by different mechanisms (a signal width and an audit probability) and can, in principle, have different-signed or different-magnitude effects on $E$, so any claim that collapses them into a single ``asymmetry level'' conflates two distinct interventions a protocol designer might separately adopt. Because $E$ is realized through a stochastic bargaining dialogue rather than a closed-form calculation, any apparent effect of $\theta$ on $E$ is entangled with which specific LLM and prompt framing generated the dialogue, and no single replication can fully separate a property of information structure from a property of the tested models. Finally, distinguishing ``$E$ scales smoothly with $\theta$'' from ``$E$ jumps at some $\theta^*$'' is a curve-shape adjudication that a coarse grid, a small per-cell sample, or a nonlinear transformation of $E$ can manufacture spuriously even when the underlying relationship is smooth, so a threshold claim requires more evidence than a single model fit.

\section{Method Overview}

At a conceptual level, the study is a single measurement pipeline with five cooperating parts. Bargaining scenarios and paired agents enter the pipeline; a ground-truth scorer turns each episode's outcome into an exploitation reading invariant to the information manipulation; a factorial sweep of the two information channels, together with a joint regression, turns the resulting series into separated per-channel effects; a claim-audit mechanism makes the verifiability channel itself manipulable and lets its effect be decomposed through a measured deception rate; a bounded replication across further models and prompts measures --- without eliminating --- how much any finding depends on the specific agents tested; and a threshold-detection gate turns the swept exploitation series into a statistically defended verdict on whether the relationship is smooth or discontinuous.

The \textbf{ground-truth exploitation scorer} computes each episode's exploitation reading arithmetically from the scenario's known reservation values and the realized outcome, never from agent beliefs or judge scores. The \textbf{disentangled axis sweep} crosses observability and verifiability factorially through two physically distinct mechanisms and fits a single regression that yields the two channels' main effects and their relative magnitude as separate, falsifiable estimands. The \textbf{claim audit and deception mediation} module extracts and checks declared claims against ground truth at a swept audit probability, producing both the verifiability manipulation itself and a deception-rate mediator. The \textbf{bounded cross-model replication} reruns the swept design across further model families and prompt framings to measure, honestly and without overclaiming, how much the finding generalizes. The \textbf{threshold detection gate} subjects the swept exploitation series to a joint statistical criterion before any claim of a discontinuity is made.

\section{Methodology}

\subsection{Ground-Truth Exploitation Scorer}
\textbf{Challenge.} Any measure of how much surplus an advantaged agent captures risks being contaminated by the very information manipulation it is supposed to quantify, if it relies on what an agent discloses, believes, or is judged to have done.

\textbf{Insight.} Anchoring the metric to evaluator-held ground truth --- reservation values and the resulting zone of possible agreement --- that is computed identically regardless of which information-structure cell produced the episode makes the metric invariant to the manipulation under study, rather than a downstream function of it.

\textbf{Method.} Each episode's exploitation reading is computed as $E = (u_a - u_a^{\text{bench}})/\mathrm{ZOPA}$, purely from the scenario manifest (true reservation values, ZOPA) and the realized outcome (deal flag, price) --- never from an LLM judge or either agent's stated beliefs. An impasse in a positive-ZOPA episode scores $E = 0$ with an accompanying efficiency-loss flag rather than being dropped from analysis.

\subsection{Disentangled Axis Sweep}
\textbf{Challenge.} Observability and verifiability are mechanistically distinct obstacles --- not knowing a counterpart's value versus not being able to check what the counterpart says --- and a single collapsed asymmetry index cannot tell a protocol designer which one to regulate, even if it correctly detects that ``asymmetry matters'' in aggregate.

\textbf{Insight.} Implementing the two channels as physically independent manipulations, rather than as two settings of one composite variable, lets a single regression specification recover two separately falsifiable estimands --- whether each channel has a real effect at all, and which channel's effect is larger --- rather than one collapsed answer.

\textbf{Method.} Observability (signal width) and verifiability (audit probability) are crossed factorially, and a mixed-effects regression of the exploitation reading on ordinal observability, verifiability, and their interaction, with random intercepts per scenario template, is fit to the resulting episode series. The fitted model yields both the marginal main effect of each channel and their standardized relative magnitude from one specification.

\subsection{Claim Audit and Deception Mediation}
\textbf{Challenge.} Making verifiability into something that can be swept, rather than a fixed yes/no property of the setting, requires a mechanism that can check a declared claim against ground truth at a controllable rate and that also produces a measurable account of \emph{why} verifiability matters, not just \emph{that} it does.

\textbf{Insight.} A graded, probabilistic audit --- rather than a binary lie-detector --- yields both a continuous independent variable (the audit probability) and a continuous mediator (the resulting deception rate) from the same mechanism, letting the verifiability channel's effect be decomposed into how much operates through actually reducing false claims.

\textbf{Method.} Factual claims about private state are extracted from the free-text claim field by a rule-assisted extractor, human-validated on a sample of claims before use; each extracted claim is checked against the scenario manifest with the swept audit probability, and detected falsehoods are publicly flagged in the transcript. The resulting per-episode deception rate (fraction of audited claims that are false) is carried into the regression as a mediator of the verifiability channel's effect.

\subsection{Bounded Cross-Model Replication}
\textbf{Challenge.} Any exploitation-vs-asymmetry relationship measured from LLM bargainers is entangled with the specific model and prompt framing used to generate the dialogue, and no finite replication can fully establish that a finding is a property of information structure rather than of the tested models.

\textbf{Insight.} Deliberately staging internal validity (one pinned model, protocol, and prompt) separately from external validity (a bounded replication across further models and prompts), rather than conflating the two, converts an otherwise-unresolved confound into a measured --- but explicitly not eliminated --- moderator.

\textbf{Method.} The swept design is rerun across several further model families and prompt framings, with model family entered as a random effect, and the sign of each channel's main effect, together with the location of any detected threshold, is compared across the tested models rather than assumed to transfer.

\subsection{Threshold Detection Gate}
\textbf{Challenge.} A coarse grid, a small per-cell sample, or a nonlinear transformation of the exploitation reading can each independently manufacture the appearance of a discontinuity even when the true underlying relationship between asymmetry and exploitation is smooth.

\textbf{Insight.} A threshold claim is trustworthy only if it survives more than one independent falsification hurdle at once --- a model-comparison test, an effect-size criterion, and replication under a denser sampling grid --- which together is what rules out a nonlinear metric or a coarse grid manufacturing a false discontinuity.

\textbf{Method.} A segmented (unknown-breakpoint) regression is compared against a smooth linear null \textbf{primarily} on the path of jointly maximal asymmetry, requiring both a formal breakpoint test and a model-fit criterion to jointly favor the segmented model before any candidate breakpoint is reported; only if this primary test is inconclusive is a \textbf{secondary} test run on each individual axis with the other held fixed. Any candidate breakpoint is then re-estimated under a denser sampling grid, and only breakpoints whose location is stable under that denser grid are reported as a threshold.

\section{Experiments}

\subsection{Experimental Setup}

\textbf{Datasets.} A synthetic reservation-value suite (uniform and log-normal draws with controlled ZOPA, generated directly by the information-structure manipulation) is used for the pilot sweep --- the only way to guarantee every tested cell has a known, controllable ground truth and a guaranteed positive-ZOPA rate. GLEE's \citep{GLEE} bargaining and negotiation game configurations supply a second, parameterized scenario family for the scale-up stage. CraigslistBargains \citep{CraigslistBargains} supplies scenario diversity --- real listings and heterogeneous price ranges --- beyond the synthetic value distributions used in the pilot, used specifically to test whether the dose-response finding and any detected threshold generalize beyond a single synthetic scenario family, rather than being reported only post hoc.

\textbf{Baselines} (sanity-check comparators, not competitors to beat --- this is a measurement study, and no research question here is framed as ``our procedure beats baseline X''). An information-symmetric anchor arm (both sides receive exact disclosure and all claims are verified) is the necessary control for every asymmetric-cell claim, and doubles as the validity check on the scorer itself. A scripted, memoryless, time-dependent concession opponent --- non-adaptive by construction --- replaces the advantaged agent to test whether any observed trend is strategic use of information or simply protocol/turn-structure arithmetic. For this control to license the causal attribution to information use, it must satisfy a pre-registered pair of criteria: the scripted opponent's own exploitation-reading trend across the observability ladder must stay within a narrow, pre-specified band of zero slope, while the LLM arm's trend must exceed the scripted arm's by a pre-specified margin per observability level; failing either criterion would mean the measured trend cannot yet be attributed to strategic use of information rather than to the bargaining protocol's own turn structure. GLEE's \citep{GLEE} own reported self-gain and fairness behavior under its information conditions is used only as a qualitative directional consistency check (does the sign of GLEE's own information-availability effects agree with this study's main-effect signs?), not a quantitative comparison, since GLEE's protocol, turn structure, and prompt scaffolding differ from this study's harness. A persuasion-gain-style measure of one-sided advantage, adapted from \citet{StrategicPersuasion}'s formal Bayesian-persuasion work, serves as a secondary, exploratory operationalization of advantage, independent of that work's own published sender-receiver results.

\textbf{Metrics.} The ground-truth exploitation index and its regression-estimated main effects and standardized relative magnitude; deal rate and efficiency loss as secondary outcomes; the audited deception rate as a mediator; and the threshold-detection module's breakpoint statistics (a formal breakpoint test's p-value, a model-fit comparison, and a bootstrap confidence interval on breakpoint location).

\textbf{Research questions:}
\begin{enumerate}[leftmargin=1.6em]
\item RQ1 --- Does the exploitation metric register zero deviation in the information-symmetric condition?
\item RQ2 --- Do observability and verifiability each have a positive causal main effect on exploitation?
\item RQ3 --- Does verifiability constrain exploitation more than raw observability?
\item RQ4 --- Does measured deception mediate the verifiability effect?
\item RQ5 --- Does the exploitation-asymmetry relationship's direction hold across model families and prompt framings?
\item RQ6 --- Does exploitation scale smoothly with asymmetry, or does it cross a genuine threshold?
\end{enumerate}

\subsubsection{RQ1: Does the exploitation metric register zero deviation in the information-symmetric condition?}
\textbf{Setup.} Within the pilot sweep's grid, the information-symmetric anchor cell (both sides receive exact disclosure, all claims verified) is analyzed separately from the asymmetric cells, with a bootstrap confidence interval computed on its mean exploitation reading.

\textbf{Expected outcome.} The symmetric anchor's mean exploitation reading is within 0.05 of zero in absolute value.

\textbf{Falsification condition.} Refuted if the bootstrap confidence interval on the symmetric anchor's mean exploitation reading excludes zero by more than 0.05 in either direction --- indicating either a defect in the scorer's computation or an unmodeled source of behavioral advantage unrelated to the information manipulation under study.

\textbf{Hypothesis (plain language).} The exploitation metric is a valid instrument: it registers no exploitation when no informational advantage has been manipulated.

\textbf{Metrics.} Bootstrap-CI mean exploitation reading in the symmetric anchor cell.

\textbf{Result.} \ph

\subsubsection{RQ2: Do observability and verifiability each have a positive causal main effect on exploitation?}
\textbf{Setup.} Pilot stage: one model pair, one synthetic scenario family, an observability-by-verifiability grid, a pre-specified number of episodes per cell. Scale-up: several backbone models (including cross-model pairs), further scenario families (GLEE, CraigslistBargains-derived), and a larger per-cell episode count. A mixed-effects model regresses the episode-level exploitation reading on ordinal observability, verifiability, and their interaction, with a random intercept per scenario template, tested by a likelihood-ratio test with Holm correction across backbones and axes at $\alpha = 0.01$.

\textbf{Sample-size rationale.} The pilot grid ($4$ observability levels $\times$ $3$ verifiability levels, 60 episodes per cell, $n = 720$) is sized for the \emph{marginal main-effect contrasts} the RQ2 test actually uses: comparing the extreme levels of one axis pools episodes across the other axis (roughly 180 episodes per marginal level), which yields approximately 90\% power to detect a standardized marginal contrast of 0.4 standard deviations at the pre-registered $\alpha = 0.01$. Individual cell-vs-cell contrasts at 60 episodes per cell are explicitly \emph{not} powered at this level (power for the same 0.4-SD contrast is roughly 0.35 at $\alpha = 0.01$) and are therefore never used as confirmatory tests on their own --- per-cell readings feed the mixed-effects fit and the pre-registered INCONCLUSIVE floor, not standalone conclusions. A pilot marginal contrast below the target pre-committedly triggers an enlarged per-cell count in the scale-up rather than a null conclusion.

\textbf{Expected outcome.} Both the observability and verifiability main effects on exploitation are positive and Holm-significant, and the maximal-asymmetry cell's mean exploitation reading exceeds the symmetric anchor's by at least 0.10.

\textbf{Falsification condition.} Refuted if either main effect is non-positive or fails to clear the Holm-corrected significance threshold, or if the maximal-asymmetry cell fails to exceed the symmetric anchor by the pre-registered 0.10 margin --- either would mean at least one of the two channels does not behave as a lever on exploitation over the tested range.

\textbf{Hypothesis (plain language).} Both losing sight of a counterpart's reservation price and being unable to verify its claims each independently increase how much surplus the advantaged agent captures.

\textbf{Metrics.} Per-cell mixed-effects main-effect estimates and confidence intervals; cell-mean exploitation reading at the maximal-asymmetry and symmetric-anchor cells.

\textbf{Result.} \ph

\subsubsection{RQ3: Does verifiability constrain exploitation more than raw observability?}
\textbf{Setup.} Standardized regression coefficients for the observability and verifiability main effects are compared directly from the model fit in RQ2, using the same specification and data. This is a pre-registered \textbf{secondary, non-gating} refinement of RQ2: it is only interpreted once RQ2 itself has resolved, and it never substitutes for RQ2. If RQ2's main-effect test fails on both axes, RQ3 is not evaluated at all (there is no verified per-axis effect for either channel to order) rather than being read as an independent ordering claim.

\textbf{Expected outcome.} Conditional on RQ2 passing, the standardized verifiability main effect's magnitude is at least as large as the standardized observability main effect's magnitude.

\textbf{Falsification condition.} Refuted if the standardized observability effect exceeds the standardized verifiability effect --- meaning raw value-observability, not claim verifiability, is the dominant lever, which would reverse which channel a protocol designer should prioritize regulating. If this test fails while RQ2 still passes, RQ2's headline dose-response wording stands unchanged and only the ordering clause is dropped, reported as ``no reliable ordering between the two axes was detected'' rather than silently asserted either way.

\textbf{Hypothesis (plain language).} Whether declared claims can be checked constrains exploitable surplus more strongly than whether the counterpart's reservation value itself is directly observable.

\textbf{Metrics.} Standardized regression coefficients for both main effects, with confidence intervals.

\textbf{Result.} \ph

\subsubsection{RQ4: Does measured deception mediate the verifiability effect?}
\textbf{Setup.} A bootstrap mediation analysis (a pre-specified number of resamples) estimates the proportion of the verifiability main effect on exploitation carried through the audited deception rate, conditional on the claim extractor clearing a human-validated reliability gate on a held-out sample of claims. Like RQ3, this is a pre-registered \textbf{secondary, non-gating} refinement of RQ2: if RQ2's verifiability main effect is not established, RQ4 is not evaluated (there is no verified verifiability effect for a mediator to explain).

\textbf{Expected outcome.} Conditional on RQ2's verifiability main effect holding, at least 30\% of that effect is mediated by the audited deception rate, with a bootstrap confidence interval on the mediated proportion excluding zero.

\textbf{Falsification condition.} Refuted if the bootstrap confidence interval on the mediated proportion includes zero, or if the point estimate falls below 30\% --- meaning the verifiability effect operates through some channel other than measured deception, such as the mere threat of an audit independent of actual lying rates. If this test fails while RQ2's verifiability effect still holds, only the mediation clause of the verifiability finding is dropped (``verifiability's effect on exploitation was not shown to route through measured deception rate; an unmeasured channel cannot be ruled out''); RQ2's status is unaffected.

\textbf{Hypothesis (plain language).} Verifiability does not just correlate with less exploitation; part of its effect is explained by agents actually lying less when their claims can be checked.

\textbf{Metrics.} Bootstrap indirect-effect estimate and confidence interval for the deception-rate mediator.

\textbf{Result.} \ph

\subsubsection{RQ5: Does the exploitation-asymmetry relationship's direction hold across model families and prompt framings?}
\textbf{Setup.} The swept design from RQ2 is rerun across several further model families and two prompt framings, with model family entered as a random effect; the sign of each main effect is compared against the pilot-stage finding within each tested family.

\textbf{Expected outcome.} The sign of both the observability and verifiability main effects agrees with the pilot-stage finding in every tested model family. This result, if obtained, bounds --- but does not by itself establish --- generality of the exploitation-asymmetry relationship beyond the specific models and prompt framings actually tested; it cannot rule out that a differently-trained or differently-prompted model class would show a different relationship.

\textbf{Falsification condition.} Refuted if the sign of either main effect disagrees with the pilot-stage finding in any tested model family --- meaning the exploitation-asymmetry relationship, or its direction, is at least partly an artifact of the specific model or prompt framing rather than a stable property of the information structure itself.

\textbf{Hypothesis (plain language).} The direction in which the two information channels affect exploitation is not an artifact of one particular model or prompt, within the bounds of what a feasible replication battery can test --- though no finite battery can establish this beyond the models and framings actually tried.

\textbf{Metrics.} Per-family standardized main-effect estimates and their signs; qualitative agreement count across tested families.

\textbf{Result.} \ph

\subsubsection{RQ6: Does exploitation scale smoothly with asymmetry, or does it cross a genuine threshold?}
\textbf{Setup.} Per Stage 1's pre-registered decision procedure, a segmented (unknown-breakpoint) regression is compared against a smooth linear null \textbf{primarily} on the path of jointly maximal asymmetry --- both channels moving together toward the extreme --- requiring both a formal breakpoint test and a model-fit criterion to jointly favor the segmented specification. Only if this primary test is inconclusive does a \textbf{secondary} test examine each individual axis with the other held fixed. Any candidate breakpoint identified in the pilot stage is re-estimated under a densified sampling grid with a larger per-cell episode count concentrated near the candidate location.

\textbf{Expected outcome.} The primary (joint-extreme) test finds a segmented model beating the smooth null, and the resulting breakpoint's location under the densified grid overlaps its pilot-stage estimate; the secondary single-axis tests are consulted only if the primary test is inconclusive, not run and reported regardless as a parallel check.

\textbf{Falsification condition.} Refuted toward ``no threshold'' if the smooth model fits at least as well as the segmented model on the primary (joint-extreme) test, and the secondary single-axis tests --- run only because the primary was inconclusive --- also fail to find a break. Refuted as an artifact --- even if a pilot-stage breakpoint initially appears --- if it fails to replicate under the densified grid, in which case the honest finding is that exploitation scales smoothly and near-additively with no regime change at any tested condition.

\textbf{Hypothesis (plain language).} Beyond a specific joint condition on both channels, exploitation stops behaving as a matter of degree and becomes qualitatively different --- or, if this is refuted, exploitation is a gradient with no cliff, and partial disclosure and partial verification each buy proportional protection.

\textbf{Metrics.} Cell-mean exploitation curve; breakpoint test statistic and model-fit comparison; bootstrap confidence interval on breakpoint location, pilot vs.\ densified grid.

\textbf{Result.} \ph

\vspace{0.5em}
\noindent\textbf{Further evaluation.} With more time, a natural extension is multi-issue and multilateral bargaining (many-to-many markets with partner choice, rather than a fixed bilateral pair) --- explicitly out of scope here because it changes both the unit of analysis and the benchmark computation (a multi-issue ZOPA and a market-level fairness reference are not a direct extension of the bilateral scorer above) and would need its own identifiability design.

\section{Related Work}

The closest prior work to this paper is GLEE, and it is worth stating the delta plainly before surveying the wider landscape: GLEE already offers a unified, parameterized family of bargaining, negotiation, and persuasion games with a configurable information-availability axis, but it treats information conditions as discrete settings to benchmark existing models under, not a continuously swept independent variable, and it has no claim-verifiability axis at all --- the exact gap this paper's audit mechanism and factorial sweep are built to close.

\textbf{Parameterized and repeated LLM-agent economic games.} GLEE \citep{GLEE} defines controllable game parameters (horizon, discounting, information availability) across bargaining, negotiation, and persuasion games with standardized outcome metrics; SOTOPIA \citep{SOTOPIA} evaluates rich multi-turn social scenarios, including negotiation, with private per-agent goals scored by a multi-dimensional LLM-judge rubric; CompeteAI \citep{CompeteAI} situates competing LLM agents in a persistent, resource-constrained market and studies cumulative advantage over many rounds. All three establish that LLM-agent behavior in configurable or repeated economic interactions is measurable, but none sweeps information structure as a continuous independent variable against a ground-truth surplus benchmark --- GLEE's conditions are discrete configurations with no verifiability axis, SOTOPIA's rubric is not surplus-anchored, and CompeteAI's endogenous feedback is precisely what prevents attributing an observed advantage to information structure rather than to its own emergent dynamics. Closer to a swept institutional manipulation, a concurrent study of GPT-based buyer/expert agents in repeated credence-goods markets \citep{InfoAsymMarkets2603} compares a verifiability institution on versus off and reports that verifiability alone provides little protection against overcharging once liability rules are absent --- a useful contrastive data point for this paper's own verifiability axis, though the comparison is a discrete categorical arm rather than a swept audit-probability ladder, the market is credence-goods quality fraud rather than bilateral price bargaining under reservation values, and no smooth-vs-threshold dose-response test is attempted. Two 2026 works come closer still and are worth conceding precisely. A study of LLM buyer-seller bargaining under partial information \citep{UsedCarSalesbots} compares complete-information, one-sided-asymmetry, and mutual-uncertainty regimes against game-theoretic solutions while measuring agents' honesty and credulity --- the nearest existing manipulation of bargaining information structure --- but its regimes are three unordered categorical arms: there is no ordinal or continuous observability ladder, no claim-verifiability mechanism of any kind, no factorial crossing of the two channels, and no test of the exploitation curve's shape. And TERMS-Bench \citep{TERMSBench} makes the bilateral price-bargaining environment itself the verifier --- the counterpart's private state is hidden from the agent but observable to the evaluator, yielding oracle-referenced surplus and optimality-gap diagnostics --- which means the evaluator-held-ground-truth scoring device this paper also uses is now shared machinery rather than a contribution; what TERMS-Bench does not do is treat the information structure as the manipulated variable at all, since its purpose is diagnosing agents, not measuring how exploitation responds to the information conditions they are placed in.

\textbf{Principal-agent and persuasion theory under private information.} A generalized principal-agent framework unifying Stackelberg games, contract design, and Bayesian persuasion \citep{LearningAgentPA} derives exploitability bounds for a principal facing a learning agent, but explicitly assumes no private information on the agent's side and evaluates regret-parameterized theory rather than measured LLM behavior. A Bayesian-persuasion formalization of debate-derived sender-receiver games \citep{StrategicPersuasion} measures a sender's persuasion gain over an uninformed receiver with a fixed prior and known utilities, but the interaction is one-directional by construction, unlike a bilateral bargain where both sides hold private information and reputational leverage. An online bilateral-trade algorithm \citep{ContextualBilateralTrade} achieves low regret on gains-from-trade under private, context-dependent valuations, but the interaction is a price-posting mechanism with one-bit feedback and no natural-language channel at all, so there is no notion of a declared claim or its verifiability. None of these three lines of work has a bidirectional, natural-language bargaining setting in which both an observability axis and a verifiability axis can be manipulated at once.

\textbf{Measuring and detecting LLM deception.} A formal, behaviorally grounded definition of AI deception in structural causal games \citep{HonestyBestPolicy} yields graphical criteria for detecting deception but no quantitative account of how much surplus deception extracts or how that scales with observability or verifiability. A text-based social-deduction sandbox \citep{AmongUs} elicits emergent deception under a hidden-role structure fixed by the game's rules, scored by Elo rating rather than a surplus-denominated outcome. A study of frontier models' harmful-query responses \citep{StrategicDishonesty} shows models can produce plausible but intentionally incorrect ``fake harmful compliance'' that exploits an evaluator's inability to verify outputs, under a fixed, binary unverifiability condition rather than a swept audit probability. A benchmark grounding inference-time privacy in contextual-integrity theory \citep{ConfAIde} tests whether models appropriately withhold information under social norms --- a compliance framing, not an adversarial one, since it does not model a counterpart strategically extracting or exploiting withheld information for material gain. Closest to this paper's verifiability channel, a real-time chain-of-thought monitor for asymmetric LLM negotiations \citep{ThinkingOutLoud} audits a seller's internal reasoning against its messages in a used-car scenario with an undisclosed defect and alerts the buyer when concealment is detected --- an oversight \emph{intervention} whose asymmetry is fixed by the scenario and whose monitoring is always-on, rather than a verifiability \emph{treatment} swept as an audit probability, and it reports walk-away and acceptance behavior rather than a surplus-denominated exploitation reading or a mediation decomposition. Together, these establish that deception and information control are measurable behaviors in LLMs, but none varies observability or verifiability parametrically or reports a surplus-denominated outcome, which is exactly what a dose-response measurement of exploitation requires.

\section{Limitations and Threats to Validity}

These are anticipated threats, stated before any experiment has run; they are to be re-examined and updated against actual outcomes once results exist.

\noindent\textbf{Construct validity --- the normative anchor.} The exploitation reading is defined against an equal-split (symmetric Nash bargaining) benchmark, but unequal splits may be normatively defensible under asymmetric protocol power (for example, a first-mover or patience advantage the protocol itself confers). Pre-registered mitigation: every hypothesis test is recomputed under a Rubinstein alternating-offers benchmark as a sensitivity analysis, and conclusions must be sign-stable across the two anchors or the discrepancy is reported rather than suppressed.

\noindent\textbf{Internal validity --- instrument reliability and thin cells.} The verifiability manipulation and the deception mediator both depend on the claim extractor, whose errors would contaminate the v-axis and the mediation analysis. Pre-registered mitigation: the extractor must clear a five-step validation sequence --- grounding the claim category in the scenario-manifest schema, wide-then-refined drafting on a held-out calibration set, validation against at least two independent human raters at a $\kappa \geq 0.8$ acceptance gate, inspection of the confusion structure for systematic direction-of-error collapse, and a freeze before any cell is scored --- before any verifiability-axis analysis; asserted reliability without this protocol is treated as absent. Separately, any grid cell whose valid-episode count falls below two-thirds of its per-cell target after gate exclusions is pre-declared INCONCLUSIVE --- never scored as null, never silently dropped --- and if this affects the joint-extreme cell, the primary threshold test itself is declared INCONCLUSIVE rather than silently replaced by the secondary per-axis tests.

\noindent\textbf{External validity --- scope and model dependence.} The design is scoped to single-issue, two-party price bargaining with computable valuations; multi-issue and multilateral markets change both the unit of analysis and the benchmark computation and are explicitly left to follow-up work. Any finding about which channel dominates, or whether a threshold exists, is established only for the tested model families and prompt framings --- the bounded cross-model replication measures this dependence but, as stated throughout, cannot remove it.

\noindent\textbf{Reproducibility --- model drift.} Findings obtained from API-served models are exposed to provider-side model drift between the pilot and scale-up stages, which could shift a measured breakpoint for reasons unrelated to information structure. Pre-registered mitigation: pinned model versions and temperature, one repeated pilot cell whose mean exploitation reading must reproduce within its bootstrap confidence interval, and one open-weights model in the replication battery.

\section{Conclusion}

This work reframes a governance question --- whether information asymmetry in agent-to-agent bargaining is a benign feature or a systemic exploitation vector --- as a measurable, decomposed quantity rather than a single existence claim. The central methodological commitment is that exploitation must be scored against ground truth the manipulation itself cannot contaminate, that its two component channels must be manipulated and estimated independently rather than folded into one index, and that any claimed threshold must survive a statistical test, an effect-size criterion, and replication under denser sampling before the word threshold is used. Whichever picture holds for current LLM bargainers --- smooth scaling or a genuine regime change --- the finding is decision-relevant: it tells a protocol designer whether partial disclosure or partial verification buys proportional protection, or whether protection is worthless short of a specific joint condition. Two limits bound this claim from the outset: the design is scoped to single-issue, two-party bargaining with computable valuations, leaving multi-issue and multilateral markets for follow-up; and any finding about which channel dominates, or whether a threshold exists, is established only for the tested model families and prompt framings, a bound the cross-model replication measures but cannot remove.

\bibliography{refs}

\vspace{1em}
\hrule
\vspace{1em}

\section*{Pre-handoff checklist}
\begin{itemize}[leftmargin=1.4em]
\item \textbf{Organization/navigation} --- narrative flows problem $\to$ method $\to$ experiments $\to$ related work without forward references to undefined terms.
\item \textbf{Terminology discipline} --- ``exploitation index,'' ``observability,'' ``verifiability,'' ``threshold,'' ``symmetric anchor'' each defined at first use in Problem Definition, used consistently thereafter; no internal stage IDs (\texttt{sp\_*}, \texttt{m\_*}, \texttt{rq\_*}) leaked into prose --- checked directly against the discussion pipeline log.
\item \textbf{Notation with worked example} --- the formal setup section defines $(b, s, r_b, r_s, \mathrm{ZOPA}, a, d, \theta, o, v, p^*, E)$ and the Problem Definition/Methodology sections consistently reuse this notation; no separate worked numeric example is included since no numbers exist yet (would require fabricating results).
\item \textbf{RQ framing matches outcomes} --- all 6 RQs are graded/quantitative (rates, effect sizes, mediated proportions, agreement counts), each answered in its own subsection with an explicit falsifiable expected outcome and a distinct falsification condition, not a binary yes/no.
\item \textbf{No typesetting artifacts} --- only explicit placeholder markers in Results, per the drafting convention this is a pre-registered spec, not a report.
\item \textbf{Artifact plan} --- the testbed harness (scenario generator, claim-audit mechanism, threshold-detection pipeline) and the self-generated transcript corpus are the natural release artifacts; not yet built, so not claimed as available.
\end{itemize}

\section*{Hand-off}
This is a draft, not a finished paper --- experiments have not run. Natural next steps: prose polish and venue-specific LaTeX formatting once a target venue is chosen, or revising an earlier stage if piloting surfaces a problem (e.g., if the maximal-asymmetry cell's exploitation reading is much smaller than expected).

\end{document}
